\section{Introduction}
  Synthetic market data generation has become an important tool in
  finance, particularly for time series such as asset returns, trading
  volume, and limit order books \cite{Potluru24, Horvath2025}. It can
  help researchers simulate fat-tail behavior, volatility clustering,
  and back-testing scenarios for prediction and hedging algorithms
  \cite{Stenger24}.

  Despite this growth, evaluation practices remain poorly
  standardized. Stenger et al. surveyed 56 studies and found little
  consensus on the metrics used or how they should be interpreted
  \cite{Stenger24}. As a result, comparisons across methods are often
  difficult and their relative strengths remain unclear.

  \subsection{A Brief Overview on Market Generators}
    Synthetic market data are multivariate time series related to
    financial assets and market microstructures, such as limit order
    books, trading volume, volatility surfaces, and asset returns
    \cite{Assefa2019, Stenger24, Horvath2025}. Researchers and
    practitioners are mainly interested in them for two reasons
    \cite{Horvath2025}: forecasting and prediction, where synthetic
    paths can support model training and evaluation, and simulating
    market scenarios, where they can reproduce rare events and
    stress-test strategies. Recent generative models such as GANs and
    VAEs, especially conditional variants, are attractive because they
    can learn conditional distributions of financial time series rather
    than only point estimators \cite{Horvath2025}.

  \subsection{Limitations in Market Generator Literature}
    Existing evaluation practices for synthetic time series remain
    fragmented, with no broadly accepted standard \cite{Wang20,
    Stenger24}. Many measures are poorly defined, inconsistently
    implemented, and hard to reproduce \cite{Stenger24}. This is
    especially problematic for generative tasks, where evaluation must
    capture more than marginal similarity and instead reflect temporal
    structure, diversity, and downstream utility \cite{Stenger24,
    Horvath2025}. Our review highlights four main limitations:

    \begin{enumerate}[label=\textbf{[L\arabic*]}]
      \item \label{limitation:metrics}
        \textbf{Inconsistent evaluation metrics:} Different studies use
        different measures, making comparisons difficult
        \cite{Stenger24}.

      \item \textbf{Lack of standardized datasets and preprocessing:}
        Benchmarking is often done on different assets, horizons, and
        preprocessing choices, which weakens comparability
        \cite{Stenger24}.

      \item \textbf{Narrow evaluation scope:} Many studies focus on
        marginal distributions and overlook temporal dependencies,
        diversity, and downstream usefulness \cite{Stenger24,
        Horvath2025}.

      \item \textbf{Reproducibility issues:} Missing implementations and
        detailed protocols make results hard to verify across studies
        \cite{Stenger24, Potluru24}.
    \end{enumerate}

  \subsection{Our Contributions}
    We develop a unified benchmark for synthetic equity-return
    generation that combines statistical and utility-based evaluation.
    Building on the taxonomy of Stenger et al. \cite{Stenger24}, the
    framework of Ang et al. \cite{Ang23}, and utility-focused work by
    Boursin et al. and Fecamp et al. \cite{Boursin22, Fecamp20}, we
    study how different metrics trade off in practice. Our key
    contributions are:

    \begin{enumerate}[label=\textbf{[C\arabic*]}]
      \item \label{contribute:1} \textbf{A consolidated evaluation
        framework:} We review and integrate measures from prior work
        to provide a more systematic benchmark for market generators
        \cite{Stenger24, Ang23, Boursin22}.

      \item \label{contribute:2} \textbf{Unified statistical and utility
        measures:} We combine fidelity metrics with downstream utility
        metrics in a single evaluation suite.

      \item \label{contribute:3} \textbf{An open benchmark framework:}
        We provide an open-source setup to support more consistent
        comparisons in future market-generator research.
    \end{enumerate}

  \subsection{Research Questions}
    Our benchmark addresses three main questions:
    \begin{enumerate}[label=\textbf{[Q\arabic*]}]
      \item How do different training sequence lengths affect generator
        performance?
      \item Are there trade-offs between fidelity, diversity, and
        downstream utility?
      \item Do modern deep-learning generators outperform classical
        parametric models in practical tasks?
    \end{enumerate}
    These questions are explored further in
    §\ref{sec:generator-framework} and the evaluation sections that
    follow.

